\chapter{Conferenz-Anzeigen}
\quellseite{12}

Die Chippewa Valley-Specialconferenz versammelt sich, f. G. m., vom 1. bis 3. Februar bei P. Lätisch in Deer Park, Wis. Arbeiten: Wie haben wir zu handeln, wenn Mitglieder unserer Aufnahme in unsere Gemeinden suchen? P. Gerike. Exegese über den ersten Brief Petri: P. Schmidtkeke. Predigtstudie: P. G. C. Meyer (P. Bubeck). Katechese: P. Steffen. Predigtvorlesen: P. Bräm. Prediger: P. Bubeck (P. Lahl). Beichtredner: P. Fuhrmann (P. Gerike). Man fahre mit der Wisconsin Central R. R. nach Eylon. Um zeitige An-, resp. Abmeldung wird dringend gebeten. Ed. Sylvester, Secr.

Zur Winterconferenz am 2. und 3. Februar in der neuen Dreieinigkeits-Kirche an 13. Ave. S. und 19. Straße zu Minneapolis, Minn., wird sofortige Anmeldung erbeten von Fr. Sievers, 1904 13th Ave. S.

Die Rock River-Pastoral- und Lehrerkonferenz versammelt sich, D. V., vom 2. bis 4. Februar in Saginaw, Wis., bei P. Hansen. Prediger: Schneider—Treu. Beichtredner: Heidelberger—Rathjen. Katechese über Frage 35: Martin—Schäfer. Behandlung eines lutherischen Kirchenliedes in der Schule: Sohn. Winke für Pastoren und Lehrer über die Worte: „Ich habe es alles Macht, aber es frommet nicht alles“: M. Müller. Exegese über Eph. 3, 13 ff.: Häuser—Wilhelm. Unterschied zwischen den Sacramenten des alten und des neuen Testaments: Theel. Man melde sich bei Zeiten! Th. Hoffmann, Secr.

Die Golfstaaten-Pastoral- und Lehrerkonferenz versammelt sich, w. G., vom 12. bis 15. Februar in P. Cramers Gemeinde zu New Orleans, La. Arbeiten: A. Schon früher angewiesen: 1. Nachweis, daß in Luthers kleinem Katechismus sämmtliche Heilslehren enthalten sind und daraus alle Irrlehren widerlegt werden können: P. Fußmann. 2. Katechese über „die bösen Engel“: Lehrer Riedel. 3. Was kann und soll der Organist thun, damit durch seinen Dienst die Gemeinde erbaut werde? (Schluß): Lehrer Rieger. 4. The Resurrection of Christ (Fortsetzung): P. Lankenau. 5. Der Zustand der Seele zwischen Tod und Auferstehung: P. Heise. 6. Christliche Schulddisciplin: Lehrer Schönhardt. B. Neue Arbeit: Die Entstehung der Concordienformel: P. Fischer. Rechtzeitige An-, resp. Abmeldung wird dringend erbeten. K. Kretzschmar, Secr.
